\chapter{Introduction}
\label{ch:intro}

This report documents the program that runs our VEX V5 competition robot. The
robot has two drive wheels, an arm, a claw, and an AI Vision sensor (a camera
that can recognise colours). The program is written in Python using the VEXcode
extension for VS Code, and it runs on the V5 Brain, the small computer on the
robot that controls the motors and reads the sensors.

A competition match has two parts. In the \emph{autonomous} period, the robot
moves on its own using only its program, so every movement must be planned and
measured in advance. In the \emph{driver control} period, a driver steers the
robot with a handheld controller, so the program's job is to turn stick
movements into smooth, predictable motor power. The program therefore has two
main halves, and most of this report explains how each half works and why it
was built that way.

The report has three aims:

\begin{enumerate}
  \item To explain the \textbf{structure} of the program: which file holds what,
    and how the files are combined before they reach the robot.
  \item To explain \textbf{what each part does}, with references to the actual
    code and worked examples of the maths behind it.
  \item To record the \textbf{timeline} of development, the \textbf{challenges}
    we faced, and how we attempted to overcome them, including the attempts
    that did not work.
\end{enumerate}

The third aim matters because the final program is only a small part of the
story. For example, the autonomous routine went through more than forty
versions in two days before we replaced it with a simpler approach. Explaining
why we made that change is as useful as explaining the final code.

\section{Sources}

Everything in this report comes from the project's own records. These are the
source code in \code{src/}, the project documentation (\code{README.md} and
\code{AGENTS.md}), and the git history. Git is a version control system: every
time we saved a meaningful change, git recorded what changed, when, and a short
message explaining why. Each saved change is called a \emph{commit} and is
identified by a short code such as \commit{eb032a0}. We quote these codes
throughout so that every claim can be checked against the history. There are
57 commits between 22 September and 26 October 2026.

Code listings in this report are read directly from the source files, so their
line numbers match the real files at the time of writing. Code that has since
been removed is reproduced from the git history and labelled as such.

\section{Outline}

\cref{ch:structure} describes the hardware and how the code is organised.
\cref{ch:movement} explains autonomous movement, including the maths that turns
a distance into a motor rotation. \cref{ch:driver} explains driver control.
\cref{ch:vision} explains how the robot recognises the cone with the AI Vision
sensor. \cref{ch:timeline} gives the day-by-day timeline, and
\cref{ch:challenges} goes through the main challenges in detail.
\cref{ch:evaluation} evaluates the program's strengths and limitations, and
\cref{ch:conclusion} concludes the report.
